\section{为什么要研究 Jacobian？}
\label{sec:jacobian-motivation}
% ============================================================================

在前五章里，我们已经分别掌握了三条线索：
一条是模曲线 $X_0(N)$ 作为紧黎曼面；
一条是权 $2$ 尖点形式 $\Sk_2(\GammaO(N))$ 作为全纯微分；
另一条是 Hecke 算子把这些微分组织成一套彼此交换的算术算子。
到这里，一个更大的问题开始浮现：这些对象怎样真正落到\textbf{椭圆曲线}和更一般的\textbf{Abel 簇（Abelian varieties）}\index{Abel簇!Abelian variety}上？如果没有新的桥梁，我们就只能分别谈“函数”“微分”“算子”，却还看不到它们如何组合成一个几何对象。

\textbf{我们遇到了什么问题？}
模形式空间 $\Sk_2(\GammaO(N))$ 是线性空间，Hecke 算子是线性算子；
但椭圆曲线、同源和有理点则生活在代数几何的世界里。
若想把“特征形式 $f$”变成“某条椭圆曲线 $E/\Q$”，光知道 $f$ 的 $q$-展开还不够，
我们需要一个几何容器，把整块权 $2$ 空间连同 Hecke 作用一起装进去。

\textbf{这个概念的核心思想是什么？}
Jacobian 簇（Jacobian variety）\index{Jacobian簇!Jacobian variety}正是这个容器。
给定一条亏格为 $g$ 的紧黎曼面 $X$，它把 $X$ 上所有全纯微分的积分周期组织成一个 $g$ 维复环面，记作 $J(X)$。
于是“曲线上的微分”被打包成“一个 Abel 簇”；Hecke 算子随后也就能从微分空间上升为 Abel 簇的自同态。

\textbf{引入之后能得到什么？}
一旦有了 $J_0(N)=J(X_0(N))$，我们就能把 Hecke 特征形式切出一个商 Abel 簇
$A_f$，这就是\textbf{模 Abel 簇（modular abelian variety）}\index{模Abel簇!modular abelian variety}。
当 $f$ 是权 $2$ 新形式且 Fourier 系数都在 $\Q$ 中时，$A_f$ 的维数就是 $1$，因而它直接是一条椭圆曲线。
这正是“模形式 $\Rightarrow$ 椭圆曲线”的桥梁；
而到了下一章的 Eichler--Shimura 理论，这座桥梁还会把两边的 $L$-函数也完全对接起来。

\begin{figure}[ht]
\centering
\begin{tikzpicture}[>=Stealth, font=\small, node distance=1.6cm]
  \node[draw, rounded corners, fill=blue!8, minimum width=4.8cm, minimum height=0.9cm] (f)
    {$f\in \Sk_2(\GammaO(N))$};
  \node[draw, rounded corners, fill=green!8, below of=f, minimum width=4.8cm, minimum height=0.9cm] (t)
    {Hecke 代数 $\mathbb T$};
  \node[draw, rounded corners, fill=orange!10, below of=t, minimum width=4.8cm, minimum height=0.9cm] (j)
    {$J_0(N)=J(X_0(N))$};
  \node[draw, rounded corners, fill=purple!10, below of=j, minimum width=4.8cm, minimum height=0.9cm] (a)
    {$A_f=J_0(N)/I_fJ_0(N)$};
  \node[draw, rounded corners, fill=red!8, below of=a, minimum width=4.8cm, minimum height=0.9cm] (e)
    {当 $[K_f:\Q]=1$ 时得到椭圆曲线 $E/\Q$};

  \draw[->, thick] (f) -- node[right] {特征值系统} (t);
  \draw[->, thick] (t) -- node[right] {对应诱导自同态} (j);
  \draw[->, thick] (j) -- node[right] {对 Hecke 理想取商} (a);
  \draw[->, thick] (a) -- node[right] {权 $2$ 且 $K_f=\Q$} (e);

  \draw[decorate, decoration={brace, mirror, amplitude=6pt}, thick]
    ($(f.west)+(-0.35,0.35)$) -- ($(j.west)+(-0.35,-0.35)$)
    node[midway, left=8pt] {解析对象与 Hecke 作用};
  \draw[decorate, decoration={brace, mirror, amplitude=6pt}, thick]
    ($(a.east)+(0.35,0.35)$) -- ($(e.east)+(0.35,-0.35)$)
    node[midway, right=8pt] {几何对象与算术应用};
\end{tikzpicture}
\caption{本章主线：从权 $2$ 模形式与 Hecke 代数出发，经由 Jacobian $J_0(N)$ 构造出与单个新形式对应的模 Abel 簇 $A_f$；在有理系数情形下，$A_f$ 就是一条椭圆曲线。}
\label{fig:ch06-motivation-bridge}
\end{figure}


从这个角度看，本章并不是突然引入一个新名词，而是在回答前面几章一直悬着的那个问题：
模曲线和 Hecke 算子已经准备好了，\emph{真正承载算术信息的几何对象到底是什么？}
答案就是 $J_0(N)$ 以及它的 Hecke 因子 $A_f$。

\begin{proposition}[前五章留下的权 $2$ 数据]
\label{prop:weight2-data-for-jacobian}
对每个正整数 $N$，都有自然同构
\[
  H^0\bigl(X_0(N),\Omega^1\bigr)\cong \Sk_2\bigl(\GammaO(N)\bigr),
\]
并且
\[
  \dim \Sk_2\bigl(\GammaO(N)\bigr)=g\bigl(X_0(N)\bigr).
\]
特别地，若 $g\bigl(X_0(N)\bigr)=1$，那么 $X_0(N)$ 上恰好有一条（差一个常数倍）非零全纯微分。
\end{proposition}

\begin{proof}
第一条同构正是\cref{thm:s2-holomorphic-differentials} 的内容：
权 $2$ 尖点形式 $f$ 与全纯微分 $f(\tau)\,d\tau$ 一一对应。
第二条等式则由\cref{cor:dim-s2-genus} 得到，因为紧黎曼面上全纯 $1$-形式空间的维数等于亏格。
最后一句只是把这个维数结论代入 $g=1$ 的情形。
\end{proof}

\begin{example}[为什么 $N=11$ 立即值得关注]
\label{ex:jacobian-motivation-11}
由\cref{ex:genus-small-levels,cor:dim-s2-genus}，
\[
  g\bigl(X_0(11)\bigr)=1,
  \qquad
  \dim \Sk_2\bigl(\GammaO(11)\bigr)=1.
\]
这意味着 $X_0(11)$ 只有一条独立的全纯微分。
对亏格 $1$ 曲线来说，这正是“它自己就是某个复环面”的信号；
本章稍后会看到，此时 Jacobian $J_0(11)$ 本身就是一条椭圆曲线。
因此 $N=11$ 是模形式真正落到椭圆曲线上的第一个非平凡例子。
\end{example}

\textbf{一句话总结。}
前五章告诉我们：$X_0(N)$ 提供了曲线，$\Sk_2(\GammaO(N))$ 提供了微分，Hecke 算子提供了可交换的算术作用。
本章要做的事，就是把这三者装进同一个对象 $J_0(N)$ 里，然后从中切出与单个特征形式对应的因子 $A_f$。
